% !TeX root = ../MQF_Manuale_Master.tex
% Capitoli/MQF_Cap_11_Programmazione_Lineare.tex

\chapter{Programmazione lineare}

\label{cap:programmazione-lineare}

\section{Obiettivi della lezione e problema guida}

\label{sec:cap11-obiettivi}

Il fallimento di Silicon Valley Bank nel marzo 2023 ha mostrato come una
strategia di bilancio possa risultare redditizia in condizioni ordinarie e,
nello stesso tempo, lasciare una banca fortemente esposta a un cambiamento
dello scenario finanziario.

La banca aveva raccolto una quota molto elevata di depositi non assicurati,
concentrati presso imprese tecnologiche e fondi di venture capital, e aveva
investito una parte rilevante delle risorse in titoli a lunga scadenza.
L'aumento dei tassi di interesse ne ridusse il valore, mentre il rallentamento
del settore tecnologico aliment\`{o} i deflussi dei depositi. Quando la banca
dovette vendere attivit\`{a} e cercare nuovo capitale, la perdita di fiducia
produsse una corsa agli sportelli di eccezionale rapidit\`{a}.

Il caso suggerisce una domanda che precede la crisi e riguarda la costruzione
del bilancio:

\begin{quote}
\emph{Come dovrebbe una banca allocare il proprio attivo per sostenere la
redditivit\`{a} senza compromettere la capacit\`{a} di fronteggiare deflussi
di liquidit\`{a} e perdite di valore, e quale costo economico comportano i
vincoli prudenziali?}
\end{quote}

La programmazione lineare permette di rappresentare questa decisione mediante
variabili decisionali, una funzione obiettivo e un sistema di vincoli. La
dualit\`{a} consente inoltre di stabilire quali requisiti limitino
effettivamente la scelta e quale valore marginale debba essere attribuito al
loro allentamento o irrigidimento.

Il percorso concettuale del capitolo \`{e}:

\[
\begin{array}{c}
\text{problema finanziario}
\longrightarrow
\text{modello lineare}
\longrightarrow
\text{regione ammissibile}
\\[5pt]
\longrightarrow
\text{soluzione ottima}
\longrightarrow
\text{problema duale}
\longrightarrow
\text{prezzi ombra}.
\end{array}
\]

Al termine del capitolo lo studente dovr\`{a} essere in grado di formulare un
problema finanziario in forma lineare, interpretarne la regione ammissibile e
i vincoli attivi, costruire il problema duale, comprendere i risultati di
dualit\`{a} debole e forte, applicare le condizioni di complementariet\`{a} e
attribuire un significato economico-finanziario alle variabili duali.

\section{Il caso Silicon Valley Bank e la decisione preventiva di bilancio}

\label{sec:cap11-caso-svb}

Silicon Valley Bank era specializzata nei servizi finanziari alle imprese
tecnologiche, alle societ\`{a} innovative e agli operatori di venture capital.
La forte espansione di questi settori tra il 2019 e il 2021 produsse una rapida
crescita dei depositi e del totale dell'attivo della banca.

Una quota molto elevata dei depositi superava il limite della garanzia
pubblica. Alla fine del 2022 circa il $94\%$ dei depositi di Silicon Valley
Bank non era assicurato. La raccolta era inoltre concentrata presso clienti
appartenenti allo stesso settore e collegati da relazioni economiche e
informative particolarmente strette.

La banca invest\`{i} una parte rilevante delle risorse raccolte in titoli
obbligazionari e mortgage-backed securities a lunga scadenza. Questi strumenti
presentavano un rischio di credito contenuto, ma erano esposti al rischio di
tasso: quando i tassi di mercato aumentarono, il valore dei titoli gi\`{a}
detenuti diminu\`{i}.

Si determin\`{o} pertanto una combinazione critica:

\[
\begin{array}{c}
\text{depositi concentrati e potenzialmente instabili}
\\
+
\\
\text{attivit\`{a} a lunga durata e sensibili ai tassi}
\\
+
\\
\text{capacit\`{a} di liquidazione immediata insufficiente}.
\end{array}
\]

Il rallentamento della raccolta di capitale nel settore tecnologico spinse
molti clienti a utilizzare i depositi per finanziare la propria operativit\`{a}.
La banca dovette quindi affrontare deflussi crescenti proprio mentre il valore
di mercato di una parte rilevante del portafoglio titoli si era ridotto.

L'8 marzo 2023 Silicon Valley Bank annunci\`{o} la vendita di circa
21 miliardi di dollari di titoli disponibili per la vendita, con una perdita
al netto delle imposte di circa 1,8 miliardi, e un progetto di aumento di
capitale. L'annuncio rese visibile la fragilit\`{a} del bilancio e contribu\`{i}
alla perdita di fiducia. Il giorno successivo i deflussi dei depositi
superarono 40 miliardi di dollari; la banca venne chiusa il 10 marzo.

Il fallimento non pu\`{o} essere ricondotto a una sola decisione di
portafoglio. Vi contribuirono carenze di governo e controllo dei rischi,
debolezze nella gestione della liquidit\`{a}, concentrazione della raccolta,
rischio di tasso e insufficiente capacit\`{a} di risposta alla corsa ai
depositi.

Il capitolo non intende ricostruire integralmente il bilancio di Silicon
Valley Bank n\'{e} valutare a posteriori tutte le decisioni assunte dalla
banca. Il caso reale viene utilizzato per formulare un problema preventivo e
stilizzato:

\begin{quote}
\emph{come allocare le risorse disponibili prima dello shock, tenendo conto
congiuntamente di redditivit\`{a}, liquidit\`{a}, perdite potenziali e
concentrazione?}
\end{quote}

Consideriamo quindi una banca che, al tempo $t_{0}$, deve distribuire il
proprio attivo tra:

\[
x_{C},
\qquad
x_{S},
\qquad
x_{B},
\qquad
x_{P},
\]

dove $x_{C}$ indica la cassa, $x_{S}$ i titoli liquidi a breve termine,
$x_{B}$ le obbligazioni a pi\`{u} lunga scadenza e $x_{P}$ i prestiti.

Le quattro classi svolgono funzioni differenti:

\begin{itemize}

\item la cassa garantisce disponibilit\`{a} immediata, ma offre una
redditivit\`{a} limitata;

\item i titoli a breve termine combinano rendimento e relativa facilit\`{a}
di monetizzazione;

\item le obbligazioni a lunga scadenza possono offrire un rendimento maggiore,
ma sono pi\`{u} sensibili a un aumento dei tassi;

\item i prestiti sostengono la funzione commerciale della banca, ma non sono
normalmente liquidabili in tempi brevi senza costi.

\end{itemize}

La decisione viene assunta prima dello stress. Il modello non descrive quindi
le operazioni di emergenza effettuate durante una corsa ai depositi, ma la
composizione dell'attivo che la banca sceglie quando dispone ancora della
possibilit\`{a} di prevenire la fragilit\`{a}.

La calibrazione numerica sar\`{a} costruita a partire dai dati e dagli ordini
di grandezza osservati nel bilancio di Silicon Valley Bank prima del
fallimento. Saranno preservate, per quanto compatibile con la struttura
semplificata del modello, la rilevanza dei depositi non assicurati, la
concentrazione del portafoglio in titoli a lunga durata, l'esposizione alle
perdite da aumento dei tassi e la limitata capacit\`{a} di fronteggiare
deflussi eccezionalmente rapidi.

Per mantenere leggibili i calcoli, gli importi potranno essere espressi in
unit\`{a} monetarie aggregate e arrotondati a zero o una cifra decimale. Le
poste contabili saranno inoltre ricondotte alle quattro classi di attivo del
modello. Tali semplificazioni non dovranno alterare il meccanismo economico
fondamentale evidenziato dal caso reale.

\section{Orizzonti temporali, ipotesi e variabili decisionali}

\label{sec:cap11-ipotesi-variabili}

La decisione viene assunta al tempo $t_{0}$, prima che si manifesti lo
scenario avverso. Il modello considera due orizzonti:

\[
\tau=t_{0}+30\mbox{ giorni},
\qquad
t_{1}=t_{0}+1\mbox{ anno}.
\]

Al tempo $\tau$ si verifica la capacit\`{a} della banca di fronteggiare
deflussi dei depositi e perdite di valore. Al tempo $t_{1}$ si valuta invece
la redditivit\`{a} annuale prodotta dalla composizione dell'attivo.

\subsection{Ipotesi del modello}

Il caso \`{e} costruito sotto le seguenti ipotesi:

\begin{enumerate}

\item la composizione dell'attivo viene scelta una sola volta al tempo $t_{0}$;

\item non sono previste decisioni di ribilanciamento successive;

\item il totale delle fonti disponibili \`{e} noto al tempo $t_{0}$;

\item rendimenti, coefficienti di liquidabilit\`{a} e perdite nello scenario
di stress sono rappresentati mediante parametri deterministici;

\item costi, rendimenti e assorbimenti sono proporzionali agli importi
investiti;

\item le quantit\`{a} decisionali possono assumere valori continui e non
negativi.

\item i parametri numerici sono ottenuti aggregando e arrotondando dati e
ordini di grandezza riconducibili al bilancio di Silicon Valley Bank; gli
aggiustamenti didattici devono preservare le principali proporzioni
economiche del caso reale;

\end{enumerate}



Le ipotesi di proporzionalit\`{a} e additivit\`{a} permettono di formulare il
problema come programma lineare. Esse costituiscono una semplificazione:
nella realt\`{a}, costi di liquidazione, rendimenti e perdite possono dipendere
dalla dimensione delle posizioni e dalle condizioni di mercato.

\subsection{Fonti disponibili}

Le passivit\`{a} iniziali sono rappresentate da:

\[
D_{I},
\qquad
D_{U},
\qquad
F_{0},
\qquad
E_{0},
\]

dove:

\begin{itemize}

\item $D_{I}$ indica i depositi assicurati;

\item $D_{U}$ indica i depositi non assicurati;

\item $F_{0}$ indica le altre fonti di finanziamento;

\item $E_{0}$ indica il capitale iniziale.

\end{itemize}

Il totale delle risorse da allocare \`{e} quindi:

\[
A_{0}
=
D_{I}+D_{U}+F_{0}+E_{0}.
\]

Le fonti sono dati del problema: la banca decide come impiegare le risorse,
non come modificarne la composizione.

\subsection{Variabili decisionali}

Le variabili decisionali sono:

\[
x=
\begin{pmatrix}
x_{C}\\
x_{S}\\
x_{B}\\
x_{P}
\end{pmatrix},
\]

dove:

\[
\begin{array}{ll}
x_{C} & \text{cassa e disponibilit\`{a} immediatamente liquide},\\
x_{S} & \text{titoli liquidi a breve termine},\\
x_{B} & \text{obbligazioni a pi\`{u} lunga scadenza},\\
x_{P} & \text{prestiti}.
\end{array}
\]

Per ciascuna classe vale il vincolo di non negativit\`{a}:

\[
x_{C},x_{S},x_{B},x_{P}\geq0.
\]

Il vettore $x$ descrive quindi la composizione preventiva dell'attivo scelta
dalla banca al tempo $t_{0}$.

\section{Costruzione del modello bancario}

\label{sec:cap11-modello-bancario}

La composizione dell'attivo deve rispettare il vincolo di bilancio:

\[
x_{C}+x_{S}+x_{B}+x_{P}=A_{0}.
\]

Indichiamo con

\[
c_{C},\qquad c_{S},\qquad c_{B},\qquad c_{P}
\]

i rendimenti annui unitari delle quattro classi di attivo. La redditivit\`{a}
complessiva \`{e} quindi:

\[
R(x)
=
c_{C}x_{C}
+
c_{S}x_{S}
+
c_{B}x_{B}
+
c_{P}x_{P}.
\]

L'obiettivo della banca consiste nel massimizzare tale quantit\`{a}:

\[
\max_{x}\;R(x).
\]

\subsection{Deflusso dei depositi e disponibilit\`{a} nello stress}

Indichiamo con

\[
\rho_{I}
\qquad\text{e}\qquad
\rho_{U}
\]

le quote dei depositi assicurati e non assicurati che vengono ritirate entro
l'orizzonte di stress. Il deflusso complessivo da fronteggiare \`{e}:

\[
\overline{u}
=
\rho_{I}D_{I}
+
\rho_{U}D_{U}.
\]

La distinzione tra le due categorie di depositi \`{e} rilevante nel caso
ispirato a Silicon Valley Bank, poich\'{e} una quota particolarmente elevata
della raccolta non era coperta dalla garanzia sui depositi e poteva quindi
reagire pi\`{u} rapidamente alla perdita di fiducia.

Non tutte le attivit\`{a} possono essere trasformate in mezzi di pagamento
nella stessa misura. Introduciamo i coefficienti:

\[
1,\qquad q_{S},\qquad q_{B},\qquad q_{P},
\]

con valori compresi tra zero e uno. Essi rappresentano la quota monetizzabile
di ciascuna classe entro l'orizzonte $\tau$.

La disponibilit\`{a} ottenibile nello scenario di stress \`{e}:

\[
Q(x)
=
x_{C}
+
q_{S}x_{S}
+
q_{B}x_{B}
+
q_{P}x_{P}.
\]

Il requisito di liquidit\`{a} impone:

\[
Q(x)\geq\overline{u}.
\]

La cassa entra con coefficiente unitario. I coefficienti delle altre attivit\`{a}
possono invece essere inferiori a uno per riflettere tempi di vendita,
riduzioni di prezzo o difficolt\`{a} di liquidazione.

\subsection{Perdite di valore e capacit\`{a} patrimoniale}

Indichiamo con

\[
m_{S},\qquad m_{B},\qquad m_{P}
\]

le perdite unitarie associate allo scenario avverso. La perdita complessiva
sull'attivo \`{e}:

\[
M(x)
=
m_{S}x_{S}
+
m_{B}x_{B}
+
m_{P}x_{P}.
\]

La cassa non compare nell'espressione, poich\'{e} nel modello non subisce
perdite di valore nello scenario considerato.

Sia $E_{\min}$ il livello minimo di capitale che la banca intende preservare.
La capacit\`{a} disponibile per assorbire perdite \`{e}:

\[
K
=
E_{0}-E_{\min}.
\]

Il vincolo patrimoniale \`{e} quindi:

\[
M(x)\leq K.
\]

I coefficienti $q_{j}$ e $m_{j}$ descrivono aspetti distinti. Il coefficiente
$q_{j}$ misura quanto dell'attivit\`{a} pu\`{o} essere monetizzato entro lo
stress; il coefficiente $m_{j}$ misura la perdita economica prodotta dallo
scenario avverso.

\subsection{Vincoli operativi e di concentrazione}

La banca mantiene almeno un livello minimo di cassa:

\[
x_{C}\geq C_{\min}.
\]

Per preservare la funzione commerciale e il finanziamento alla clientela,
impone inoltre:

\[
x_{P}\geq P_{\min}.
\]

La concentrazione nelle obbligazioni a lunga scadenza viene limitata mediante:

\[
x_{B}\leq\alpha_{B}A_{0},
\]

dove $\alpha_{B}$ rappresenta la quota massima dell'attivo destinabile a tale
classe.

\subsection{Programma lineare completo}

Il problema decisionale assume quindi la forma:

\[
\begin{array}{ll}
\max
&
c_{C}x_{C}
+
c_{S}x_{S}
+
c_{B}x_{B}
+
c_{P}x_{P}
\\[4pt]
\mbox{soggetto a}
&
x_{C}+x_{S}+x_{B}+x_{P}=A_{0},
\\
&
x_{C}+q_{S}x_{S}+q_{B}x_{B}+q_{P}x_{P}
\geq\overline{u},
\\
&
m_{S}x_{S}+m_{B}x_{B}+m_{P}x_{P}
\leq K,
\\
&
x_{C}\geq C_{\min},
\\
&
x_{P}\geq P_{\min},
\\
&
x_{B}\leq\alpha_{B}A_{0},
\\
&
x_{C},x_{S},x_{B},x_{P}\geq0.
\end{array}
\]

Il modello traduce il problema finanziario in una struttura composta da una
funzione obiettivo e da un insieme di decisioni ammissibili. La soluzione
ottima sar\`{a} la composizione dell'attivo che massimizza la redditivit\`{a}
annua rispettando tutti i requisiti assegnati.